\chapter[American Civilisations]{How American Civilisations Developed in Different Environments}
\section{A Knotted Cord}
Pick up an object now lying in a museum drawer in Lima, Peru. Even through glass, a photograph reveals its form.

A horizontal main cord supports dozens of thinner pendants like hanging plaits. Brown, red, yellow, and blue strings carry knots, large and small, solitary near the top or crowded below. It might feel like a woven ornament or a lost instrument's component.

This is a quipu, ``knot'' in Quechua. Five centuries ago, in the vast Andean Inca state, it was an archival document. Trained officials whose title meant keepers of knots read it with their fingers: a red cord for village maize taxes, a yellow one for stored cloth, positions for units, tens, and hundreds, colours for goods, perhaps twist directions for further meanings. Some scholars suggest non-numerical records too, including genealogy, events, or songs. That remains unsettled because scarcely anyone alive can fully read them.

Pause over the discomfort. We equate information with writing: letters in ledgers, paper archives, history books. Yet these cords kept accounts and governed a country thousands of kilometres long without a single written character.

Did the Inca therefore have writing or not? Behind this apparent word game lies our larger question. Without Eurasia's standard equipment---iron, alphabetic writing, horses, cattle, wheeled vehicles---can a society be complex? Do our measures of complexity and advancement assess the thing itself or its resemblance to one particular ruler?

The knotted cord begins our inspection of that ruler.

\section{The Great City on the Lake}
Now enter a scene.

Time: a November morning in 1519. Place: today's Mexico City, then a great highland lake with an island capital, Tenochtitlan, belonging to the Mexica, commonly called Aztecs today.

You are fifteen or sixteen, paddling from a lakeside village before dawn. Cool grey-blue water surrounds dozens and hundreds of canoes like market traffic. Some carry maize and beans, others baskets of tomatoes and chillies, gobbling turkeys, or artisans like you holding woven mats and pottery to try their luck.

The city arrives first by smell: toasted tortillas, water plants, and thousands of cooking fires. Then market voices swell across the water. Finally, white houses spread before you, temple platforms rise, and ruler-straight causeways reach from shore with antlike pedestrians.

You enter canals beside floating gardens: not natural land but generations of dredged mud, enclosed by fencing and rooted with willows into long, narrow artificial fields producing several annual vegetable crops. A gardener straightens to greet you, the ground gently swaying beneath him.

At Tlatelolco market, you tie up and enter a dizzying crowd. Spanish soldiers later recalled tens of thousands daily and more goods than they could remember: gold dust in goose-quill tubes, glittering quetzal-feather cloaks, rows of shaving-sharp obsidian blades, and individually counted cacao beans used as small change for turkeys or firewood. Robed judges settle short-weight disputes; tax collectors register incoming goods and take shares in kind rather than asking coin prices, for there are no coins.

You might notice absences: horses, cattle, or large domestic beasts carrying or pulling loads; heavy goods move on human backs. No iron or ringing smithies, only volcanic-glass obsidian for the best cutting. But you probably notice as little as fish notice water. You know only a city of perhaps one or two hundred thousand, larger and cleaner than almost any Spanish city, with freshwater aqueducts and organised rubbish collection.

You do not know that this very month, hundreds of bearded men riding unfamiliar great animals are entering along a causeway. Their expressions match yours on first seeing their horses.

\section{A Report Card with Missing Subjects}
State the conflict before conclusions.

Older histories often judge civilisation through a checklist: writing, metallurgy especially iron, wheeled transport, large draught animals, cities, and states. It was not deliberately designed against the Americas but generalised Eurasian experience: farming and surplus supported non-farmers, then cities, writing, bronze, and empires followed in Mesopotamia, Egypt, India, and China.

Score the Americas around 1500 and gaps glare. Iron: absent; tools and weapons were largely stone, wood, and obsidian, with some late Andean and western Mexican bronze chiefly for ritual rather than farming. Draught animals: absent; llamas and alpacas carried modest loads, not people, ploughs, or chariots. Wheeled transport: absent. Alphabetic writing: absent in most places.

Stop there and backwardness appears obvious.

But do not close the report card. Other facts occupy the same continent:

Maya cities had flourished for over a millennium in southern Mexico and Guatemala's tropical forests. Tikal, Palenque, and Copán had pyramid temples, ballcourts, and inscribed history. Priest-scholars combined 260-day sacred and 365-day solar calendars into a fifty-two-year round, with Long Count dates fixed on a millennia-long timeline. They calculated Venus's cycle with errors measured in hours and independently developed a zero sign, achieved in very few places, India another. Their elaborate hieroglyphic script recorded kings, wars, and astronomy in bark-paper folding books.

Central Mexico's Teotihuacan grew between the first and sixth centuries into a city perhaps a hundred thousand strong, when Roman splendour neared its end. Today's Avenue of the Dead leads to the Pyramid of the Sun. Even its name, roughly City of the Gods, came from Aztecs viewing ruins nearly a millennium later. We do not know the builders' self-name or language. A nameless society built one of ancient America's largest cities.

In less than a hundred years during the fifteenth century, the Inca expanded thousands of kilometres along the Andes, governed millions, built tens of thousands of road kilometres, and terraced steep slopes. All without written administrative documents, using our opening cords.

In today's Illinois, beside the Mississippi, Cahokia gathered ten or twenty thousand people around 1050, perhaps more, and built over a hundred earthworks. Monks Mound's base rivals Egypt's Great Pyramid, making it ancient North America's largest structure. Residents assembled in plazas and used timber circles to observe sunrise and schedule farming. Later colonisers saw strange grassy mounds and refused Native American authorship, inventing alternatives because their checklist offered no such box.

The conflict is clear. Missing equipment coexists with great cities, precise astronomy, domesticated maize, and governed empires. Must one side be wrong, or did we measure the wrong thing?

Keep a provisional judgement; you will revise it as we continue.

\section{Amazement, Flames, and Living People}
Different positions reveal different Americas. Hear several voices.

\textbf{First voice: conquerors' amazement.}

Soldier Bernal Díaz del Castillo saw Tenochtitlan with Cortés's force in 1519. In his later True History of the Conquest of New Spain, he remembered waterborne towns and straight causeways resembling enchanted castles in chivalric romances; some soldiers wondered whether they dreamed. This was no sentimental poet but a veteran familiar with European cities. His ruler measured something beyond prior experience.

Amazement did not prevent destruction two years later. Díaz also recorded fighting, burning, and massacre. The same hands admired and destroyed, a recurring chill in history: recognising civilisation's value does not prevent humans destroying it.

\textbf{Second voice: testimony from the ruins.}

Decades later, friar Bernardino de Sahagún invited Mexica noble descendants and young people to speak and write their history, beliefs, and lives in Nahuatl, producing the bilingual Florentine Codex. Its war account records burning temples, looted treasure, and displaced survivors from their side. Later collections such as Broken Spears preserve this uncommon testimony of defeat. We hear it not through conquerors' kindness but because a small group stubbornly recorded it.

\textbf{Third voice: allies' choices.}

Often omitted, the capture was not a few hundred Spaniards defeating hundreds of thousands of Aztecs. Tens of thousands of Indigenous allies formed the bulk of their force, notably Tlaxcalans long pressured, blockaded, and drawn into captive-taking ritual Flower Wars by the Aztec Triple Alliance. Cortés appeared neither a descending god nor simply an invader, but a useful powerful ally against an old order. America already had rivalries, alliances, and calculations; it was no silent continental backdrop awaiting discovery.

\textbf{Fourth voice: a mixed-heritage descendant's memories.}

Garcilaso de la Vega, related to the final Inca rulers through his royal mother and son of a Spanish officer, later lived in Europe. His Royal Commentaries drew on childhood memories from Cusco to describe roads, stores, law, and festivals, insistently telling readers what the lost country was really like. Idealised recollection still restored the Inca from pagan ruins to an order deserving serious description.

Now perform a crucial exercise: \textbf{state the backwardness argument fully}.

Fairness requires its strongest case before dismissing Eurocentric arrogance. First, collision produced real outcomes: Tenochtitlan fell in 1521, Atahualpa was captured in 1532, and steel blades, armour, horses, and gunpowder conferred genuine asymmetry. Second, Old World smallpox arrived before Spaniards among populations without immunity, linked to Eurasia's long livestock and disease history. Third, recording limitations were real: Maya writing survived only in limited regions before conquest, while quipu struggled with complex narrative, leaving law and history vulnerable in human memory to fire or broken transmission.

These arguments fail not in their facts but in conclusion. They establish \textbf{different environments producing different solutions, with collision's costs borne chiefly by one side}, not the absence of complex civilisation. Evidence supports the former and refutes the latter. Separate Tenochtitlan's fall as event, steel and disease as causes, and backwardness as judgement. Judgement requires a ruler whose origins we must inspect.

Correct another common error: mysterious Maya disappearance, undeciphered writing, and people evaporating overnight. Popular twentieth-century books loved this largely false story. Many southern Maya cities declined in the eighth and ninth centuries amid drought, war, and royal crises, but their people relocated rather than vanished. Millions now speak dozens of Mayan languages in Guatemala and southern Mexico and sustain ancestral calendars and practices. Descendants of a vanished civilisation scroll social media today. The disappearance belongs to our map's blind spot, not the Maya.

\section[Thinking Bridge: Solving Problems]{Thinking Bridge: Ask How, Not Whether They Had It}
Turn this into a portable tool.

Asking whether America had writing, wheels, or iron judges by an \textbf{object checklist}, assuming universal solutions. Civilisations actually solve \textbf{problems}: accounts, distant messages, steep-slope food production, heavy transport. Objects vary with environment; problems remain.

Replace it with a \textbf{needs checklist}. Ask not whether they possessed X but how they solved the problem. Practise fully on the Andes:

\noindent
\begin{tabularx}{\linewidth}{llX}
\toprule
Problem & \shortstack[l]{Common Eurasian\\ solution} & Inca solution \\
\midrule
\shortstack[l]{Tax and stock\\ records} & \shortstack[l]{Written ledgers, rods,\\ abacuses} & Quipu: numbers encoded in knots, colours, positions \\
\shortstack[l]{Distant orders and\\ news} & \shortstack[l]{Post horses and\\ carriages} & Relay runners on mountain roads, staffed stations day and night \\
\shortstack[l]{Farming steep\\ slopes} & \shortstack[l]{Great-river plain\\ irrigation} & Mountainside terraces conserving soil and water across elevations \\
Famine relief & \shortstack[l]{Public granaries, as\\ in China} & Nationwide stores collecting in good years and distributing in disasters \\
Taxation & Money and grain levies & No currency: direct labour for roads, farming, and military service \\
Crossing gorges & \shortstack[l]{Stone bridges and\\ ferries} & Woven grass-rope suspension bridges renewed annually \\
\bottomrule
\end{tabularx}

\smallskip
None of the right-hand solutions is an inferior imitation. Each answers Andean conditions independently. Roads served walkers and llama caravans; horses could not run stepped mountain paths anyway. Labour and storehouse redistribution did not require currency. Knots recorded the state's essential numbers quickly and reliably without fearing rainforest damp.

Consider two sharper examples.

\textbf{The wheel counterexample.} Claims that Americans never invented wheels are instructively wrong: Mexican wheeled pottery animals show the concept existed without wheeled transport. Vehicles need suitable roads and large draught animals; plateaux, gorges, and steep slopes, with llamas carrying only modest loads, offered poor returns. Invention's fate depends not only on conceiving it but on conditions making it worthwhile. When someone condemns a society for lacking X, ask which problem X solves there and whether another solution existed.

\textbf{Writing's shared starting point.} Writing did not begin for poetry and history, but accounting. Early Mesopotamian tablets list beer and sheep, the same business as quipu. Humanity's earliest writing impulse was inventory, not lyric expression. From the needs perspective, quipu is not undeveloped writing but another mature accounting answer.

Use the tool beyond America. Assess schools by solved student problems, not swimming pools; assess yourself through your needs rather than another person's possessions. Objects deceive; needs do not.

\section{Maize People and Knotted Government}
Read two primary materials from continents unknown to one another, addressing two sides of one matter.

The K'iche' Maya Popol Vuh passed orally through generations before sixteenth-century nobles wrote it in Latin letters and eighteenth-century friar Francisco Ximénez copied and preserved it. Gods first made unsuccessful mud people, dissolved by rain, then wooden people lacking spirit and remembrance of their makers. Finally yellow and white maize formed flesh and blood: these humans spoke, gave thanks, and looked to the stars.

Remember the equation: \textbf{people are made of maize}. For a maize-dependent society organising its year around cultivation, this grounds a worldview, not merely fairy-tale rhetoric. Bodies come from fields; tending fields tends oneself. Calendrical precision and solemn sacrifice rest on this foundation. The Chinese saying that people are iron and food steel expresses a diluted intuition: food defines identity, not just nutrition.

The second passage, shorter, comes from the Book of Changes' Great Commentary, Part II:

\begin{quote}
In high antiquity they governed with knotted cords; later sages replaced these with written tallies.
\end{quote}
Eight characters: ancient people managed affairs with knots, later replaced by writing. A writer living amid established literacy remembered, or inherited memory of, Chinese knotted government. Cords and accounts were not an Andean eccentricity but a shared past. Eurasia travelled from cords to script; the Andes stayed longer and developed cords further into imperial administration. Two routes, one starting point.

Remember the limits. The Popol Vuh reveals K'iche' understandings, not all Indigenous American beliefs: Maya, Mexica, Inca, and Cahokian traditions differed as European ones did. The Great Commentary's eight characters establish a remembered Chinese cord age, not how actual ancient knots were tied. Evidence comes in pieces; imagination needs restraint.

\section{Three Explanations for Another Route}
With facts and tools, ask why American development differed from Eurasia. These are causal explanations, with distinct strengths and blind spots.

\textbf{Explanation one: the environment dealt the cards.}

Jared Diamond's Guns, Germs, and Steel argues that continental differences came chiefly from available resources, not different human qualities. Most American large mammals, including wild horses, camels, and mastodons, disappeared around the last ice age's end, leaving few domestic candidates: llamas, alpacas, turkeys, and dogs. Without cattle and horses came no ploughs, chariots, or post horses, and fewer livestock-linked disease sources, unlike the Old World history producing devastating smallpox. Eurasia's east--west breadth allowed crops to travel similar latitudes and climates, as wheat reached China from western Asia. The Americas' north--south length and Panama's rainforest forced maize through changing climates towards eastern North America or the Andes. Llamas never crossed north through the isthmus.

The cards genuinely differed; denying that ignores evidence. But environmental destiny goes too far. Conditions set difficulty, not answers. Without horses, Mexica built lake cities, Inca mountain roads, Cahokians great mounds. People chose these differing solutions.

\textbf{Explanation two: diffusion and pathways.}

Look at specific chains. Technologies need accumulating exchange: Eurasian smiths, farmers, and sailors converted inventions elsewhere into new starting points, from Chinese cast iron to Indian numerals and Arab navigation. American networks also spread maize, pottery, and ballgames, with Andean metallurgy and Mexican jade traditions. But forests, deserts, mountains, and lower population and density fragmented exchange, shortening and slowing chains. Paths had inertia: excellent obsidian reduced pressure for metal; adequate quipu reduced urgency for writing. Reasonable steps accumulated into a route far from Eurasia's.

This restores concrete choices missing from environmental accounts, but describes how more than why. Overused, contingency obscures real constraints.

\textbf{Explanation three: the ruler is faulty.}

Inspect the question instead. Who made writing, iron, wheels, and draught animals mandatory? Eurasian, especially European, experience became universal marking criteria. Failure is predictable when another answer sheet supplies the standard. Ask which indicators fairly compare different routes: population, farmers needed to support one noble, astronomical precision, road length, art? Change rulers and rankings change.

But rejecting a faulty ruler does not forbid comparison. Steel and stone, smallpox and immunity made real battlefield differences and contributed to millions of deaths and conquest. Correcting unfair evaluation must not erase consequences.

Environment describes the hand, pathways its play, rulers the referee. Not wholly compatible, each reveals something true. Mature thinking seats all three together.

\section[Further Challenge: Complexity Is Complex]{Further Challenge: Complexity Is Itself Complex}
We have repeatedly used complex civilisation without inspecting it. Anthropologists and archaeologists have spent over a century doing so.

Mid-twentieth-century anthropologist Elman Service classified increasing organisational scale as \textbf{bands}, kin-linked groups of dozens; \textbf{tribes}, hundreds or thousands linked through clans and villages; \textbf{chiefdoms}, thousands or tens of thousands under hereditary rank; and \textbf{states}, with bureaucracy, law, armies, and taxation. This converted complexity into measurable categories. Most classify Cahokia as a chiefdom, with debate over statehood; Inca state organisation is unquestioned; Maya cities shared culture while governing separately, like Greek city-states.

Beware the hidden \textbf{staircase illusion}.

The sequence invites assumptions that every society evolves towards states, others merely stalled, and greater complexity means superiority and goodness. The classification authorises neither. Archaeology shows chiefdoms forming then dispersing, and peoples familiar with neighbouring states choosing other arrangements. Complexity costs: states build roads and stores but also demand labour and wage war; hierarchy organises projects and inequality. Cahokia declined after two brilliant centuries for uncertain environmental or political reasons, but dispersed descendants did not fail; they carried mound-building traditions into other North American societies.

The trap has a painful form in Aztec captive sacrifice and religious warfare. Understanding is not defence, explanation not absolution. Flower Wars combined faith, training, and noble politics; Mexica believed sacrifice sustained the cosmos. Complete logic does not erase the concrete person on the altar. Complexity and cruelty can coexist with lake cities and Venus calculations. Refusing one stamp of advanced or barbaric respects facts' actual complexity.

Use complexity as an assessment report, not a prize. It describes how many people coordinate, how many levels exist, how long the accumulated chains are, not whether people live well or how they deserve treatment. A medical report does not establish character.

\section{Half the Americas on Your Table}
Return to today. The examination of our rulers continues in another room.

Think of dinner: maize, potatoes, tomatoes, chillies, sweet potatoes, pumpkins, peanuts, cacao, and more were domesticated in the Americas. Maize leads global food output; potatoes helped Eurasian population growth; tomato-free Mediterranean or chilli-free Sichuan cuisine now seems unimaginable. The civilisation marked deficient helped feed the Old World over five centuries. Post-Columbian species exchange was among history's most consequential trades, with food itself an American export.

Other survivals lie beyond meals. Mexico City stands on Tenochtitlan's ruins; drained-lake temple foundations emerge beside the cathedral. Mexico's emblem retains the founding legend's eagle with snake on cactus. Quechua is official in Peru and Bolivia, spoken daily by tens of millions; Andean farmers still use some Inca-era terraces. Quipu remains incompletely deciphered as mathematicians and anthropologists compare thousands of knot patterns in databases, making imperial archives a modern research frontier.

Cahokia offers another reminder. Many mounds lay silent beneath grass for centuries, some cut away for roads and houses, until serious protection in the later twentieth century and World Heritage listing in 1982. A former metropolis needed a change in newcomers' vision to become visible again. Unseen does not mean absent.

The ruler still measures people: one score for a whole pupil, one accent for knowledge, GDP or rankings for a city's worth, or inability to do X for an entire group's inferiority. Each repeats a particular environment's object list as universal standards. Standards need not disappear. Before scoring, ask whose experience produced the list, whether it measures objects or needs, and whether alternative solutions exist. Scores may remain; arrogance may diminish.

\section{Your Question: Who Made the Ruler?}
Return to the museum cord.

Over five centuries it has been called ornament, primitive toy, proof of absent writing, and now an incompletely deciphered information system. The cord stayed the same; eyes and rulers changed.

Answer the final question with reasons:

\textbf{If two civilisations never met---no 1492, conquest, or smallpox, each continuing on its continent until today---would asking which was more advanced still make sense?}

If yes, choose a fair ruler: population, lifespan, art, astronomy, hunger frequency, or something else. If no, face a challenge: when actual encounters imposed real costs, could we still describe asymmetry? Where is the boundary between recognising difference and declaring inferiority?

Turn it towards yourself. Exams, admissions, income, and housing will repeatedly score your life on others' lists. If you made your own, what would it contain? Knotted cords can archive an empire; can your own report card contain a real self others fail to read?

There is no standard answer. Next time someone hands you a ready-made checklist, take another look. That glance is the chapter's entire purpose.
